\documentclass[11pt]{article}

\usepackage[margin=1in]{geometry}
\usepackage{amsmath,amssymb,amsthm,mathtools,microtype}
\usepackage[T1]{fontenc}
\usepackage[hidelinks]{hyperref}

\newtheorem{theorem}{Theorem}[section]
\newtheorem{lemma}[theorem]{Lemma}
\newtheorem{corollary}[theorem]{Corollary}
\newtheorem{remark}[theorem]{Remark}

\newcommand{\C}{\mathbb C}
\newcommand{\Z}{\mathbb Z}
\newcommand{\N}{\mathbb N}
\newcommand{\E}{\mathbb E}
\newcommand{\CT}{\operatorname{CT}}
\newcommand{\ord}{\operatorname{ord}}
\newcommand{\conv}{\operatorname{conv}}
\newcommand{\Supp}{\operatorname{Supp}}

\title{A Filtered Frobenius Proof of the\\Two-Variable Gaussian Moments Conjecture}
\author{Roy van Rijn\\{\small Independent researcher}}
\date{23 July 2026\\Revised 17 November 2026}
\hypersetup{
  pdftitle={A Filtered Frobenius Proof of the Two-Variable Gaussian Moments Conjecture},
  pdfauthor={Roy van Rijn},
  pdfsubject={The Gaussian Moments Conjecture in two real variables},
  pdfkeywords={Gaussian Moments Conjecture, Gaussian moments, constant terms,
    Laurent polynomials, constant-term theorem}
}

\begin{document}
\maketitle

\begin{abstract}
The Gaussian Moments Conjecture asks whether a polynomial whose positive
powers all have Gaussian integral zero must eventually remain zero after
multiplication by any fixed polynomial.  We give an alternative proof in
two real Gaussian variables of the support classification and sharp
cutoff also obtained by Wilson.  The proof uses circular coordinates to separate
rotation from radius.  The angular part becomes a constant-term problem,
while the radial part becomes the simple functional \(U^j\mapsto j!\).
An exposed lower face of the resulting Newton diagram, the constant-term
theorem of Duistermaat and van der Kallen, and reduction modulo a prime then
isolate a nonzero lowest term.  This forces every polynomial with vanishing
pure moments to have rotational weights of one strict sign, from which the
mixed-moment conclusion follows with the sharp cutoff \(m>\deg Q\).
Together with the
counterexamples of Long in every dimension at least three, this determines
the exact dimensional range of the conjecture.
\end{abstract}

\section{Introduction}

Let \(X_1,\ldots,X_d\) be independent standard real Gaussian random
variables.  The Gaussian Moments Conjecture, introduced by Derksen, van den
Essen, and Zhao~\cite{DVZ}, is the following statement.

\medskip
\noindent
\textbf{Gaussian Moments Conjecture \(\operatorname{GMC}(d)\).}
If \(P\in\C[X_1,\ldots,X_d]\) satisfies
\[
  \E(P^m)=0\qquad\text{for every }m\geq 1,
\]
then, for every \(Q\in\C[X_1,\ldots,X_d]\),
\[
  \E(QP^m)=0\qquad\text{for all sufficiently large }m.
\]

\medskip
The conjecture was motivated in part by its connection with the Jacobian
Conjecture~\cite{DVZ}.  Long constructed explicit counterexamples to
\(\operatorname{GMC}(d)\) for every \(d\geq3\)~\cite{Long}.  The
two-variable conclusion is also contained in the contemporaneous results
of Wilson~\cite{Wilson} and Long~\cite{LongTorus}.

\begin{theorem}\label{thm:main}
The Gaussian Moments Conjecture holds in two real variables.  More
precisely, its premise implies
\[
  \E(QP^m)=0\qquad(m>\deg Q)
\]
for every nonzero polynomial \(Q\).  This uniform cutoff is optimal.
\end{theorem}

Wilson proves the same one-sided support classification and
the cutoff \(m>\deg Q\)~\cite[Theorem~1]{Wilson}; his proof also uses Newton faces
and prime isolation.  The argument below packages the arithmetic step as
a filtered Frobenius lemma.

In two
dimensions, rotation has a single integer weight and there is a single
radial variable.  Thus every polynomial can be displayed schematically as
\[
  P=\sum_k \underbrace{T^k}_{\text{rotation}}\,
                  \underbrace{B_k(U)}_{\text{radius}}.
\]
Gaussian integration first takes the constant term in \(T\), then replaces
\(U^j\) by \(j!\).  If both positive and negative rotational weights occur,
the lowest radial orders have a supporting line above weight zero.  The
polynomial carried by that lower face has a nonzero constant term in some
power, by a theorem of Duistermaat and van der Kallen~\cite{DK}.  Raising
that power to a suitable prime and reducing modulo the prime makes the
lowest term impossible to cancel after the factorial substitution.  This
contradicts the assumed moment vanishing.

Long's contemporaneous manuscript~\cite[Theorems~1.1--1.2]{LongTorus}
proves a stronger factorially weighted constant-term theorem with one radial
variable and an arbitrary-rank torus.  Its rank-one specialization gives
the same Gaussian classification, including the orthogonal-nullcone
description.  It also characterizes eventual pure vanishing and detects
nonvanishing at every sufficiently large prime dilation of one fixed power.

The contribution here is an alternative filtered Frobenius argument and
its formal interfaces; the classification and cutoff are known conclusions.
The bivariate Lean supplement starts from an ordinary polynomial in \(X,Y\)
and verifies the passage through circular coordinates to mixed moments,
subject to the two mathematical inputs described in Section~7.

\section{Circular coordinates}

Let \(X,Y\) be independent standard real Gaussian variables.  Set
\[
  Z=\frac{X+iY}{\sqrt2},\qquad
  W=\frac{X-iY}{\sqrt2}.
\]
Wick's formula gives
\begin{equation}\label{eq:wick}
  \E(Z^aW^b)=
  \begin{cases}
    a!,&a=b,\\
    0,&a\neq b.
  \end{cases}
\end{equation}
Introduce a radial variable \(U\) and a Laurent variable \(T\) through the
injective substitution
\[
  Z\longmapsto T,\qquad W\longmapsto UT^{-1}.
\]
Thus \(Z^aW^b\) becomes \(T^{a-b}U^b\).  Every polynomial \(P(Z,W)\)
therefore has a unique expression
\begin{equation}\label{eq:weight-decomposition}
  P=\sum_{k\in S}T^kB_k(U),
\end{equation}
where \(S\subset\Z\) is finite and every \(B_k\in\C[U]\) is nonzero.  The
integer \(k\) is the rotational weight.

Define the factorial functional
\[
  \mathcal L:\C[U]\longrightarrow\C,\qquad
  \mathcal L(U^j)=j!.
\]
Equation~\eqref{eq:wick} is exactly the identity
\begin{equation}\label{eq:expectation-factorization}
  \E(F)=\mathcal L\bigl(\CT_T F\bigr)
  \qquad(F\in\C[Z,W]).
\end{equation}
This is the feature special to two variables: there is one angular
constant-term operation and one radial factorial.

For example, \(P=Z+W\) becomes \(T+UT^{-1}\).  Its first angular
constant term is zero, but
\[
  \CT_T(P^2)=2U,\qquad \E(P^2)=2.
\]
The two opposite weights balance at the lowest radial degree.  For a
general polynomial, cancellations can hide this balance in low powers;
the next two sections explain why they cannot hide it in every power.

\section{The exposed lower face}

For each \(k\in S\), write
\[
  \nu_k=\ord_U B_k,\qquad
  b_k=[U^{\nu_k}]B_k\neq0.
\]
The points \((k,\nu_k)\) record the smallest radial degree at each
rotational weight.  We now extract the lower face visible from weight zero.

\begin{lemma}[Lower face]\label{lem:face}
Suppose \(0\in\conv(S)\).  There are rational numbers \(\rho,\theta\), a
subset
\[
  S_0=\{k\in S:\nu_k=\rho+\theta k\},
\]
and an integer \(r\geq1\) with the following properties:
\[
  \nu_k\geq\rho+\theta k\quad(k\in S),\qquad
  0\in\conv(S_0),
\]
and, for
\[
  P_0(T)=\sum_{k\in S_0}b_kT^k,\qquad
  F_m(U)=\CT_T(P^m),
\]
there are \(n\in\Z_{\geq0}\) and \(c\in\C^\times\) such that
\begin{align}
  &\text{every term of }F_m\text{ has degree at least }\rho m,
       \label{eq:global-bound}\\
  &n=\rho r,\qquad F_r(U)=cU^n+\text{terms of degree greater than }n.
       \label{eq:leading}
\end{align}
\end{lemma}

\begin{proof}
Let
\[
  C=\conv\{(k,\nu_k):k\in S\}\subset\mathbb R^2.
\]
Since \(0\in\conv(S)\), the polygon \(C\) meets the vertical line \(k=0\).
Let \((0,\rho)\) be the lowest point of this intersection.  The lower
boundary of \(C\) admits a supporting line through this point; write it as
\[
  \nu=\rho+\theta k.
\]
Thus
\[
  \nu_k\geq\rho+\theta k.
\]
The vertices of \(C\) are integral, so its lower edges have rational
slopes.  If \((0,\rho)\) is a vertex, choose a rational slope between the
slopes of its adjacent lower edges; at an endpoint, choose a rational
finite slope beyond the single adjacent slope.  If \(C\) is a single point,
take \(\theta=0\).  Thus \(\rho,\theta\in\mathbb Q\).

The supporting line cuts out a face of \(C\), and this face contains
\((0,\rho)\).  Every point of the face is a convex combination of its
vertices, which are among the points indexed by
\[
  S_0=\{k\in S:\nu_k=\rho+\theta k\}.
\]
Taking first coordinates in such a convex combination for \((0,\rho)\)
gives \(0\in\conv(S_0)\).

We use the constant-term theorem of Duistermaat and van der
Kallen~\cite[Theorem~2 and Remark~3]{DK}: if every positive power of a complex Laurent polynomial
has zero constant term, then zero does not lie in the convex hull of its
support.  Applied contrapositively to \(P_0\), it gives an \(r\geq1\) for
which
\[
  c:=\CT_T(P_0^r)\neq0.
\]

Consider a weight-zero product of \(m\) terms selected from
\eqref{eq:weight-decomposition}.  If the selected weights are
\(k_1,\ldots,k_m\), then its radial degree is at least
\[
  \sum_{a=1}^m\nu_{k_a}
  \geq \rho m+\theta\sum_{a=1}^m k_a
  =\rho m.
\]
This proves~\eqref{eq:global-bound}.  When \(m=r\), equality can occur only
by taking the lowest coefficient from weights on \(S_0\), and the sum of
all equality terms is \(c\).  Hence the lowest term of \(F_r\) is
\(cU^{\rho r}\).  Its exponent is necessarily a nonnegative integer; call
it \(n\).
\end{proof}

\section{Filtered Frobenius isolation}

For any characteristic-zero ring \(R\), write
\[
  \mathcal L_R:R[U]\longrightarrow R,\qquad
  \mathcal L_R(U^j)=j!.
\]
The arithmetic heart of the proof is independent of circular coordinates.

The normalization must take place before reduction: dividing by a
factorial in characteristic \(p\) would usually be meaningless.  Instead,
the radial order bound makes every factorial quotient an integer, and
these integer weights can then be reduced.

\begin{lemma}[Filtered Frobenius isolation]
\label{lem:filtered-frobenius}
Let \(R\) be a characteristic-zero domain, let \(\kappa\) be a field of
characteristic \(p>0\), and let \(R\to\kappa\), \(a\mapsto\bar a\), be a
ring homomorphism.  Suppose \(F_m\in R[U]\) satisfy, for some
\(\rho\in\mathbb Q\), \(r\geq1\), \(n\in\mathbb Z_{\geq0}\), and
\(c\in R\setminus\{0\}\),
\[
  \ord_UF_m\geq\rho m,\qquad
  n=\rho r,\qquad
  F_r=cU^n+\text{terms of degree greater than }n.
\]
If
\[
  \bar c\neq0,\qquad
  \overline{F_{rp}}=\overline{F_r}^{\,p},
\]
then the values \(\mathcal L_R(F_m)\) cannot all vanish.
\end{lemma}

\begin{proof}
Suppose they do.  The order bound at \(m=rp\), together with
\(n=\rho r\), gives
\[
  \ord_UF_{rp}\geq \rho rp=np.
\]
Thus \((np)!\) is a common factor of all factorial weights in
\(\mathcal L_R(F_{rp})=0\).  It is nonzero in the characteristic-zero
domain \(R\), so it may be cancelled before reduction.

Write \(F_r=cU^n+U^{n+1}G(U)\).  After reduction, the Frobenius hypothesis
gives
\[
  \overline{F_{rp}}
  =\overline{F_r}^{\,p}
  =\bar c^{\,p}U^{np}+U^{(n+1)p}\overline{G}(U)^p.
\]
Writing the second term as
\(\sum_{j\geq n+1}a_jU^{jp}\), the cancelled moment identity reduces to
\[
  0=\bar c^{\,p}
    +\sum_{j\geq n+1}a_j\frac{(jp)!}{(np)!}.
\]
Every factorial quotient in the sum contains the factor \((n+1)p\), hence
vanishes in characteristic \(p\).  Therefore \(\bar c^{\,p}=0\), contrary
to \(\bar c\neq0\).
\end{proof}

\section{Proof of the theorem}

\begin{proof}[Proof of Theorem~\ref{thm:main}]
Write \(P\) as in~\eqref{eq:weight-decomposition} and suppose that
\(\E(P^m)=0\) for every \(m\geq1\).  By
\eqref{eq:expectation-factorization}, this says
\[
  \mathcal L(F_m)=0\qquad(m\geq1).
\]
If \(P=0\), every positive mixed power vanishes, so assume \(P\neq0\).
If \(0\in\conv(S)\), Lemma~\ref{lem:face} produces a nonzero lowest
coefficient \(c\).  Let \(R\subset\C\) be the \(\Z\)-algebra generated by
the coefficients of \(P\) and by \(c^{-1}\).  Choose a maximal ideal
\(\mathfrak m\subset R\).  Since \(\Z\) is Jacobson and \(R\) is of finite type over \(\Z\),
\(\mathfrak m\cap\Z=(p)\) for a prime \(p\), and
\(\kappa=R/\mathfrak m\) is a finite extension of \(\mathbb F_p\)
\cite[Tags 00GB and 00GC]{StacksProject}.  The image of \(c\) is nonzero because
\(c\) is a unit in \(R\).  Moreover, in characteristic \(p\),
\[
  \overline{F_{rp}}
  =\CT_T\bigl((\bar P^r)^p\bigr)
  =\bigl(\CT_T(\bar P^r)\bigr)^p
  =\overline{F_r}^{\,p},
\]
since Frobenius sends \(T^i\) to \(T^{pi}\).  Lemma
\ref{lem:filtered-frobenius}, with the bounds from Lemma~\ref{lem:face},
now gives a contradiction.  Therefore
\[
  0\notin\conv(S).
\]
Since \(S\subset\Z\), either every weight of \(P\) is strictly positive or
every weight is strictly negative.

Now fix a polynomial \(Q\).  If \(Q=0\), there is nothing to prove, so let
\(S_Q\) be its nonempty finite set of rotational weights.  Suppose first
that all weights of \(P\) are positive, and put
\[
  a=\min S>0,\qquad b=\min S_Q,\qquad
  N=\max\left\{0,\left\lfloor\frac{-b}{a}\right\rfloor+1\right\}.
\]
Every weight of \(QP^m\) is at least \(b+ma>0\) for \(m\geq N\).

If instead all weights of \(P\) are negative, put
\[
  a=-\max S>0,\qquad b=\max S_Q,\qquad
  N=\max\left\{0,\left\lfloor\frac{b}{a}\right\rfloor+1\right\}.
\]
Then every weight of \(QP^m\) is at most \(b-ma<0\) for \(m\geq N\).
Thus in either case \(\CT_T(QP^m)=0\) for every \(m\geq N\).
Equation~\eqref{eq:expectation-factorization} then gives
\(\E(QP^m)=0\).  To obtain the stated uniform bound, observe that every
weight of a degree-\(q\) polynomial lies in \([-q,q]\), while \(a\geq1\).
Thus \(m>q=\deg Q\) suffices for either sign.

For sharpness, take \(P=Z\) and \(Q=W^q\), where \(q\geq1\).  Then
\(\E(P^m)=0\) for every \(m\geq1\), but
\(\E(QP^q)=\E(Z^qW^q)=q!\ne0\).
\end{proof}

\begin{corollary}\label{cor:classification}
The Gaussian Moments Conjecture holds in dimension \(d\) exactly when
\(d\leq2\).
\end{corollary}

\begin{proof}
The case \(d=1\) is known~\cite[Proposition~4.2]{DVZ}.
Theorem~\ref{thm:main} proves the case \(d=2\).
Long's examples disprove the conjecture for every \(d\geq3\)~\cite{Long}.
\end{proof}

\section{Why one radial variable matters}

Long's multitorus theorem~\cite{LongTorus} shows that the angular weight
lattice may have arbitrary rank while the same isolation works.  The
essential restriction is the single radial order and its factorial \(j!\).
In ordinary higher-dimensional Gaussian space there are several
independent radial directions with product-factorial weights, so the
one-dimensional isolation step no longer applies.
Corollary~\ref{cor:classification} shows that this is not merely a
technical gap: counterexamples already occur in dimension three.

\section{Formal verification}

The versioned companion Lean~4 supplement~\cite{FormalGMC2} starts
with an ordinary bivariate polynomial and the circular Wick expectation.
Its exported theorem asserts eventual mixed vanishing from vanishing of
all positive pure moments.  The checked interfaces include the circular
substitution, Wick formula, rational supporting face, lowest radial
coefficient, specialization to a finite residue field, normalized factorial
cancellation, Frobenius commutation and prime isolation.  The package
contains no \texttt{sorry} declarations.

Two mathematical inputs remain named axioms: the Duistermaat--van der
Kallen constant-term theorem and good finite-characteristic specialization
of finite-type $\mathbb Z$-domains.  The recorded endpoint audit lists
these two inputs besides Lean's foundational choice, quotient and
propositional-extensionality axioms.  Thus the formal theorem is conditional
on those interfaces.  The explicit cutoff and optimality are proved in
the proof of Theorem~\ref{thm:main} and are not separately exported.  The supplement supplies the
pinned toolchain and dependencies, build target, exact endpoint and
instructions for reproducing its axiom report.

\paragraph{Acknowledgment.}
The author thanks Christopher D. Long for making his more general
factorially weighted torus theorem publicly available~\cite{LongTorus}.

\paragraph{AI assistance and review.}
ChatGPT (OpenAI) assisted with mathematical arguments, literature checks,
exposition, proof review and verification code.

\begingroup
\small
\bibliographystyle{plain}
\bibliography{references}
\endgroup

\end{document}
